\section{Consideraciones finales}
\label{sec:consideraciones}

La auditoría documental mostró una prevalencia importante de fallas inferenciales críticas en artículos cuantitativos de revistas sudamericanas de acceso abierto. La comparación directa con estudios previos debe hacerse con cautela porque las unidades de análisis, los criterios de codificación y los dominios temáticos no son idénticos. Sin embargo, nuestros hallazgos son consistentes con la literatura que documenta deficiencias persistentes de reporte y de interpretación estadística en múltiples disciplinas \citep{ioannidis2005,darrigo2024,schmidt2025}.

Inéditamente, la prevalencia estricta de \textit{Fallas Fuertes} alcanzó el 39,0\,\% de los artículos empíricos con inferencia encuestados, mientras que una proporción todavía mayor (49,7\,\%) incurrió en procedimientos similares escudados tras una advertencia de limitación. Sumadas, casi 9 de cada 10 manuscritos exhiben incongruencias metodológicas-inferenciales, lo que representa una prevalencia notablemente superior a los rangos reportados en auditorías previas de calidad en otras regiones. Aunque la heterogeneidad de criterios hace improcedente la comparación unívoca, los hallazgos de \citet{schmidt2025} sobre el mal uso del valor $p$ sugieren que este orden de magnitud no es atípico bajo definiciones estrictas de \textit{falla inferencial}.

Un antecedente regional útil es el estudio de \citet{peralta2023}, el cual evaluó artículos de revistas y reportó cumplimientos altos de criterios APA. Aunque ese estudio se enfoca en calidad general de formato, resulta ilustrativo porque sugiere que una estructura formal aceptable no garantiza ninguna correspondencia con el rigor inferencial: un artículo puede estar bien formateado y aun así presentar problemas severos de muestreo, extrapolación e interpretación.

A nivel internacional, la literatura sobre lineamientos editoriales también ha mostrado vacíos recurrentes en el reporte de elementos esenciales \citep{schulz2010}. Si bien esto se conoce en ensayos clínicos, nuestro estudio refuerza la idea de que omisiones de diseño y precisión permeabilizan todo el abanico disciplinario, afectando desde Ciencias de la Salud hasta Ciencias Exactas, Naturales y Educación.

Una contribución central del estudio es metodológica. Al transitar hacia la recuperación y clasificación documental asistida por Inteligencia Artificial y confirmación humana, incrementos exponenciales en los volúmenes de revisión ---llegando a más de 600 artículos con clasificación integral--- se hacen posibles, reduciendo la dependencia sobre muestras pequeñas de conveniencia. Asimismo, la definición explícita de falla inferencial estricta como un problema rotundo (\textit{Falla Fuerte}) elimina ambigüedades. El protocolo diferenció penalizaciones por desconocimiento total frente a debilidades declaradas, lo cual añade densidad al debate sobre si basta con justificar una debilidad metodológica para continuar esgrimiendo pruebas inferenciales con aparente normalidad.

La fase de simulación permitió separar con claridad dos planos. En el plano probabilístico, \SRS\ y la estratificación proporcional produjeron coberturas próximas a la nominal, junto con sesgos reducidos y \RMSE\ acotado. En el plano no probabilístico, la selección por conveniencia generó sesgo estructural y volvió improcedente la interpretación de cobertura poblacional. Este resultado es coherente con la literatura metodológica. Las fórmulas de tamaño muestral y los \IC\ clásicos suponen un mecanismo probabilístico de selección, por lo que su traslado automático a muestras de conveniencia carece de justificación \citep{noordzij2010}. Del mismo modo, las advertencias sobre el uso de poder post-hoc, inferencia desanclada del diseño y reexpresiones numéricas de la incertidumbre sin base probabilística apuntan en la misma dirección \citep{hoenig2001,thomas1997,reinhart2015,navarrete2019}. Nuestra simulación no pretende resolver un debate teórico abierto, sino mostrar cuantitativamente que, bajo el mecanismo de conveniencia modelado, el sesgo persiste aunque se apliquen procedimientos de remuestreo.

Por esa razón, el bootstrap se interpretó solo como diagnóstico de variabilidad condicional a la muestra. Este punto es importante porque en la práctica aplicada suele utilizarse como sustituto informal de la aleatorización. Nuestros resultados muestran que, en los escenarios evaluados, el remuestreo no corrige la incoherencia originaria entre diseño de selección y validez inferencial.

La pertinencia regional del estudio radica en que el ecosistema sudamericano de acceso abierto combina diversidad disciplinaria, multilingüismo y modelos editoriales con recursos heterogéneos \citep{beigel2024}. En ese contexto, las exigencias estadísticas suelen ser menos visibles que las exigencias formales de estilo o estructura. Nuestros hallazgos sugieren que la mejora editorial no depende solo de reforzar normas generales de presentación, sino de exigir información mínima sobre población objetivo, marco muestral, mecanismo de selección, pesos, varianza y alcance de la inferencia. Los hallazgos del estudio sugieren que las revistas no deberían limitar la evaluación metodológica a la presencia de pruebas estadísticas o a la estructura formal del manuscrito. La revisión editorial debe verificar coherencia entre población objetivo, mecanismo de selección, estimador, método de varianza e interpretación. Esta recomendación es consistente con el espíritu de STROBE, CONSORT y SAMPL, pero aquí se operacionaliza en una herramienta breve orientada a la coherencia inferencial \citep{vonelm2007,schulz2010,lang2015}. El Apéndice~\ref{app:checklist} presenta un checklist mínimo que traduce estos principios en preguntas concretas y verificables, susceptible de integrarse a las instrucciones para autores, a formularios de revisión por pares o a listas internas de control editorial.

El estudio presenta, no obstante, varias limitaciones que deben explicitarse. En primer lugar, el marco editorial puede presentar subcobertura si existen revistas sudamericanas pertinentes fuera de las fuentes utilizadas para construir el universo. En segundo término, la clasificación del tipo de muestreo depende del reporte disponible en los artículos, por lo que es posible cierta misclasificación cuando los autores describen de forma incompleta el mecanismo real de selección. En tercer lugar, algunos subgrupos son pequeños, en especial el idioma inglés, por lo que los resultados estratificados deben interpretarse de manera descriptiva. Finalmente, la fase de simulación representa un mecanismo específico de conveniencia; otras formas de selección no probabilística podrían producir magnitudes distintas de sesgo y \RMSE.

Entre las líneas futuras de investigación, resulta particularmente relevante ampliar el marco editorial a otras bases y actualizar la auditoría de forma longitudinal; estimar prevalencias por categoría específica de falla, además del indicador compuesto; incorporar no respuesta y mecanismos de selección más complejos en la simulación; y evaluar el efecto de políticas editoriales o checklists estadísticos sobre la prevalencia de fallas a lo largo del tiempo.

En relación con auditorías previas, este estudio aporta un diagnóstico con diseño probabilístico, ponderación e \IC\ basados en varianza de diseño. En relación con la literatura sobre muestreo no probabilístico, aporta una demostración cuantitativa de la pérdida de validez inferencial bajo conveniencia en los escenarios evaluados. En relación con las guías editoriales, aporta una traducción operativa de sus principios hacia un checklist centrado en población, selección, precisión, varianza e interpretación.

El patrón observado no se reduce a errores aislados de redacción, sino que expresa una desconexión más profunda entre población objetivo, mecanismo de selección, cálculo de tamaño muestral, estimación de varianza e interpretación final. En términos prácticos, el aporte del estudio es doble. Por un lado, ofrece un diagnóstico con base de diseño sobre la prevalencia de fallas inferenciales críticas. Por otro, propone una herramienta editorial breve para fortalecer la coherencia entre diseño, análisis e interpretación. Ambos aportes convergen en una misma conclusión: mejorar la calidad de la investigación cuantitativa requiere pasar de un énfasis exclusivo en significancia y formato a un estándar explícito de transparencia metodológica y precisión inferencial.
